% !TEX program = pdflatex
\documentclass[journal]{IEEEtran}

% ──────────────────────────────────────────────────────────────────────────────
%  Packages
% ──────────────────────────────────────────────────────────────────────────────
\usepackage[T1]{fontenc}
\usepackage{amsmath,amssymb,amsfonts,amsthm}
\usepackage{algorithmic}
\usepackage[ruled,linesnumbered]{algorithm2e}
\usepackage{graphicx}
\usepackage{textcomp}
\usepackage{xcolor}
\newcommand{\rev}[1]{#1}% colors disabled for submission
\newcommand{\revB}[1]{#1}% colors disabled for submission
\newcommand{\greenrev}[1]{#1}% colors disabled for submission
\newcommand{\redrev}[1]{\textcolor{red}{#1}}% red: implementation corrections
\usepackage{booktabs}
\usepackage{multirow}
\usepackage{hyperref}
\usepackage{cleveref}
\crefname{lemma}{Lemma}{Lemmas}\Crefname{lemma}{Lemma}{Lemmas}% blue: ensure Lemma X renders correctly (shared counter with theorem)
\crefname{corollary}{Corollary}{Corollaries}\Crefname{corollary}{Corollary}{Corollaries}
\usepackage{bm}
\usepackage{mathtools}
\usepackage{subcaption}
\usepackage{siunitx}
\usepackage{balance}

% Theorem environments
\newtheorem{theorem}{Theorem}
\newtheorem{lemma}[theorem]{Lemma}
\newtheorem{corollary}[theorem]{Corollary}
\newtheorem{definition}[theorem]{Definition}
\newtheorem{remark}{Remark}
\newtheorem{assumption}{Assumption}
\newtheorem{proposition}{Proposition}

% Notation shortcuts
\newcommand{\Cfree}{\mathcal{C}_{\mathrm{free}}}
\newcommand{\Cspace}{\mathcal{C}}
\newcommand{\IR}{I_{\!R}}
\newcommand{\IE}{I_{\!E}}
\newcommand{\dR}{d_{\!R}}
\newcommand{\SPD}{\mathbb{S}_{++}^{d}}
\newcommand{\muR}{\mu_{\!R}}
\newcommand{\Vol}{\mathrm{Vol}}
\newcommand{\diag}{\mathrm{diag}}
\DeclareMathOperator*{\argmin}{arg\,min}

% IEEE specific
\IEEEoverridecommandlockouts

\begin{document}

% ══════════════════════════════════════════════════════════════════════════════
%  TITLE
% ══════════════════════════════════════════════════════════════════════════════
\title{RIT*: Riemannian Informed Trees for Cost-Adaptive Optimal Motion Planning}

% \author{
%     Muhayy~Ud~Din,
%     Ahmed~Nadar,
%     Jan~Rosell,
%     and~Irfan~Hussain
%     \thanks{Manuscript submitted on 10th April 2026.}
%     \thanks{This work was supported by Khalifa University under Award
% Nos. RC1-2018-KUCARS-8474000136,  MBZIRC-8434000194, and KU-DFL-8475000016. }
%     \thanks{Muhayy~Ud~Din,
%     Ahmed~Nadar,
% , and Irfan~Hussain are with Khalifa University Center for Autonomous Robotic Systems (KUCARS), Khalifa University, United Arab Emirates (e-mail: muhayyuddin.ahmed@ku.ac.ae, 
% irfan.hussain@ku.ac.ae).}%
%     \thanks{Jan~Rosell,is with Institute of Industrial and Control Engineering, Universitat Politècnica de Catalunya, Barcelona, Spain(e-mail: jan.rosell@upc.edu).}
%     \thanks{Code, data, and scripts to reproduce all reported results are provided as supplementary material accompanying this submission. The full implementation will be released publicly upon acceptance.}
% }
\markboth{IEEE Robotics and Automation Letters}{Author~et~al.: RIT*: Riemannian Informed Trees}

\maketitle

% ══════════════════════════════════════════════════════════════════════════════
%  ABSTRACT  (~150 words)
% ══════════════════════════════════════════════════════════════════════════════
\begin{abstract}

%Finding optimal collision-free paths in high-dimensional configuration
%spaces remains challenging, particularly under anisotropic and
%spatially varying costs such as those induced by joint inertia or
%obstacle proximity.
%Existing informed sampling-based planners restrict search to a subset
%of the space, but define this region using Euclidean distance, leading
%to inefficient sampling when the \rev{underlying cost structure} is non-Euclidean.

We present Riemannian Informed Trees (RIT*), a planning framework
that replaces Euclidean primitives in batch-informed search with their
Riemannian counterparts.
RIT* constructs a tighter, cost-consistent informed set, performs
a nearest-neighbour search under an anisotropic distance metric, and evaluates
edge costs efficiently via a cascading scheme.
\rev{We further introduce a Collision-Adaptive Metric Refinement (CARM),}
\rev{which learns an obstacle-proximity cost field online from collision feedback,}
reducing the reliance on prior metric design in practical settings.
Experiments across environments from 2-D to 14-D show that RIT* is competitive in low-dimensional and spatially constant-metric settings and produces substantially lower-cost solutions when the metric varies spatially in high-dimensional configuration spaces. Performance gains scale with anisotropy and dimension, reaching up to 13.0\% improvement over BIT* in the 3-D anisotropic benchmark, up to 24.0\% in 6-DOF manipulation, and 24.8-63.5\% in a 14-DOF bimanual planning problem, where Euclidean-informed baselines degrade.
\end{abstract}

\begin{IEEEkeywords}
Motion planning, Riemannian geometry, asymptotic optimality,
informed sampling, online metric learning.
\end{IEEEkeywords}
\vspace{-4mm}
% ══════════════════════════════════════════════════════════════════════════════
%  OVERVIEW FIGURE
% ══════════════════════════════════════════════════════════════════════════════



% ══════════════════════════════════════════════════════════════════════════════
%  I. INTRODUCTION
% ══════════════════════════════════════════════════════════════════════════════
\section{Introduction}\label{sec:intro}
Sampling-based motion planning is a fundamental technique in robotics, which
enables collision-free path computation in high-dimensional configuration spaces~\cite{LaValle2006}. The development of asymptotically optimal (AO) planners, including Rapidly-exploring Random Trees* (RRT*)~\cite{Karaman2011} and Probabilistic Roadmap Method* (PRM*), demonstrated that random sampling methods can asymptotically converge to an optimal path. Later work improved convergence by employing \emph{informed sampling}: after an initial solution with some cost is obtained, all subsequent samples are then taken solely from the \emph{informed set}, which corresponds to a prolate hyperspheroid in Euclidean space~\cite{Gammell2018}. This idea is the basis for Informed~RRT*~\cite{Gammell2018}, BIT*~\cite{Gammell2020}, AIT*, EIT*~\cite{Strub2022}, and APT*~\cite{APT2025}. 


These planners show promising performance; however, they select this set using Euclidean distance, which does not adequately reflect the true cost structure in anisotropic environments.
In addition, existing approaches typically rely on a predetermined cost structure rather than learning it from planning experience, which reduces their effectiveness when the cost landscape is unknown or varies with the environment.

\begin{figure}[t]
\centering
\includegraphics[width=0.9\columnwidth]{figures/abstract.png}
\caption{Overview of RIT*.
\textbf{(a)}~BIT* samples from the Euclidean informed set~$\IE$, which ignores anisotropic cost structure.
\textbf{(b)}~RIT* uses a \rev{cost-aligned} Riemannian informed set aligned with the underlying cost structure, while the CARM module (red shading) learns a cost field online from collision feedback, guiding paths away from obstacles without prior knowledge.}
\vspace{-5mm}
\label{fig:abstract}
\end{figure}

In this work, we address both limitations with RIT*. We replace
Euclidean components with a Riemannian formulation, where a spatially
varying metric tensor~$G(x)$ captures the \rev{underlying cost structure}.
\rev{We further introduce a Collision-Adaptive Metric Refinement (CARM),}
\rev{which learns a collision-induced conformal cost field while planning from collision feedback and}
\greenrev{adapts the planner accordingly.}

\textbf{Contributions:}
The key contributions are:
\begin{enumerate}

\item \textit{Riemannian informed planning:}
RIT* replaces Euclidean components with Riemannian ones. \rev{The informed set better aligns with the underlying cost structure, nearest neighbours use an anisotropic metric, and edge costs are computed via a cascading scheme to reduce computation.} These changes focus sampling on relevant regions and improve efficiency.
\item \textit{Online proximity-aware metric refinement via CARM:}
We introduce CARM, which learns an obstacle-proximity cost field from collision feedback and applies it as a conformal scaling of a base metric. It requires no prior obstacle knowledge: given a base metric (e.g., joint-space inertia), CARM refines it with proximity-aware structure; without one, it defaults to Euclidean and learns from scratch.

\item \textit{Experimental validation:}
We evaluate RIT* from \mbox{2-D} and \mbox{3-D} to high-dimensional tasks (6-DOF, 14-DOF), including a real \revB{UR10e} arm, and provide a comprehensive benchmark against asymptotically optimal planners.
\end{enumerate}
\vspace{-3mm}
\section{Related Work}\label{sec:related}

We review prior work on informed sampling-based planning, the use of
Riemannian geometry in motion planning, and obstacle-aware cost design.

\textit{Euclidean informed sampling planners:}
Sampling-based planners have been improved by restricting search to an
informed subset of the space. Gammell~et~al.~\cite{Gammell2018}
introduced direct sampling of a prolate hyperspheroid to accelerate
RRT*. BIT*~\cite{Gammell2020} combines batch sampling with informed sets
and lazy edge evaluation, while AIT* and EIT*~\cite{Strub2022} extend
this with asymmetric bidirectional search. APT*~\cite{APT2025} shapes
nearest-neighbour search using virtual forces and adapts batch size, and
FIT*~\cite{FIT2024} improves early convergence through batch adaptation.
Despite these advances, all define the informed set, nearest-neighbour
search, and edge cost in \emph{Euclidean} distance.

\textit{Riemannian geometry in motion planning.}
Riemannian geometry naturally models anisotropic costs, but its use is
largely limited to local or optimisation methods. CHOMP~\cite{Ratliff2009}
and STOMP~\cite{Kalakrishnan2011stomp} optimise trajectories under such costs, but
are not sampling-based and lack AO guarantees. In sampling-based
planning, Kim~et~al.~\cite{Kim2016} use Riemannian metrics for local
steering, and Kyaw and Kelly~\cite{Kyaw2026} propose a Riemannian RRT*
using retractions and natural gradients. However, these methods are
incremental, do not construct informed sets, and do not address global
sampling. Bi-AM-RRT*~\cite{Zhang2023biamrrt} similarly biases
nearest-neighbour selection via a diffusion-based metric, but does not
integrate informed sampling or batch processing.

\textit{Obstacle-aware Riemannian metrics.}
Several works incorporate obstacle information into metric design but
typically rely on prior knowledge or demonstrations.
Klein~et~al.~\cite{Klein2023} design barrier-based metrics requiring full
obstacle geometry, while RMPflow~\cite{Cheng2018rmpflow} composes
task-space metrics for reactive control. Beik-Mohammadi
et~al.~\cite{BeikMohammadi2023} learn Riemannian manifolds from
demonstrations for reactive motion generation. These approaches are not
integrated with sampling-based planning and require information beyond
collision checking.

\textbf{Key distinction:}
Prior Riemannian approaches use the metric primarily for \emph{local} operations, such as steering, trajectory optimisation, or reactive control. In contrast, RIT* integrates Riemannian geometry into the
\emph{global} informed sampling pipeline. It replaces the informed set, nearest-neighbour search, edge cost evaluation, and rewiring operations of
a batch-informed planner with their Riemannian counterparts. Additionally, RIT* incorporates CARM, which \rev{learns a conformal scaling online} from collision feedback, enabling adaptive and cost-aware planning
without prior environment knowledge.
%══════════════════════════════════════════════════════════════════════════════
%  II. PROBLEM FORMULATION
% ══════════════════════════════════════════════════════════════════════════════
\section{Problem Formulation}\label{sec:problem}

We consider optimal motion planning in a configuration space where the
cost of motion is anisotropic and may vary across the space. Such costs
are naturally described by a Riemannian metric, which assigns a
direction-dependent traversal cost at each configuration.

Let $\Cspace \subseteq \mathbb{R}^d$ denote a $d$-dimensional
configuration space, and let $\Cfree \subseteq \Cspace$ be the subset of
collision-free configurations. The space is equipped with a smooth,
positive-definite metric tensor field
$G : \Cspace \to \SPD$, where $\SPD$ denotes the set of symmetric
positive-definite matrices. For any configuration $x \in \Cspace$, the
matrix $G(x)$ defines the local cost of moving through $x$, so that
different directions can have different traversal costs.

Under this metric, the cost of a path $\sigma : [0,1] \to \Cspace$ is
its Riemannian arc length,
\begin{equation}\label{eq:cost}
  c(\sigma) = \int_0^1 \!\sqrt{\dot{\sigma}(t)^\top G\bigl(\sigma(t)\bigr)\, \dot{\sigma}(t)}\; dt,
\end{equation}
where $\dot{\sigma}(t)$ is the tangent vector of the path. The resulting
\emph{geodesic distance} $\dR(x,y)$ is defined as the minimal cost of any path between the two
configurations $x$ and $y$, i.e.,
$\dR(x,y) = \inf_\sigma c(\sigma)$ taken over all paths $\sigma$ that connect
$x$ and $y$. In the special case where $G(x) = I_d$, this corresponds to the usual Euclidean
distance.

Given a start state $x_s \in \Cfree$ and a goal state $x_g \in \Cfree$,
our objective is to find a collision-free path of minimum cost:
\begin{equation}\label{eq:opt}
  \sigma^* = \argmin_{\sigma \,:\, x_s \to x_g,\; \sigma \subset \Cfree} c(\sigma),
\end{equation}
where $\sigma$ is any feasible path that remains in the free space. A
planner is \emph{asymptotically optimal} if the cost of its returned
solution, denoted by $c_n$, converges almost surely to the optimal cost
$c^*$ as the number of samples $n \to \infty$.

After a feasible solution with cost $c_{\mathrm{best}}$ is found, the
search can be restricted to configurations that may still generate a
better path. In Euclidean informed planning, this region is the
\emph{informed set}
\begin{equation}\label{eq:IE}
  \IE = \bigl\{ x \in \Cspace : \|x_s - x\|_2 + \|x - x_g\|_2 \le c_{\mathrm{best}} \bigr\},
\end{equation}
which contains all points that could lie on a path shorter than the
current best solution under Euclidean distance.

For anisotropic costs, however, Euclidean distance does not accurately
reflect the true cost-to-come and cost-to-go. We therefore define the
\emph{Riemannian informed set} as
\begin{equation}\label{eq:IR}
  \IR = \bigl\{ x \in \Cspace : \dR(x_s, x) + \dR(x, x_g) \le c_{\mathrm{best}} \bigr\}.
\end{equation}
This set contains configurations that can still improve the current solution under the Riemannian cost. \rev{The relationship between $\mathcal{I}_R$ and $\mathcal{I}_E$ depends on the metric spectrum. When $G \succeq I_d$, $d_R(x,y)\ge \|x-y\|_2$, resulting in $\mathcal{I}_R \subseteq \mathcal{I}_E$. More generally, under the bi-Lipschitz bounds of Sec.~\ref{sec:analysis}, the two informed sets are equivalent up to constant-factor distortion, while $\mathcal{I}_R$ remains aligned with the underlying cost structure.}

The goal of this work is to design an AO informed planner that operates natively under the Riemannian cost: sampling from $I_R$, selecting neighbours and evaluating edges according to the geodesic distance $d_R$, and, when no base metric is available or when the available metric does not encode obstacle proximity, refining the metric online from collision feedback gathered during planning.
% ══════════════════════════════════════════════════════════════════════════════
%  III. METHOD
% ══════════════════════════════════════════════════════════════════════════════


\section{Method}\label{sec:method}
%\subsection{Overview}\label{sec:overview}
Batch Informed Trees~\cite{Gammell2020} is an asymptotically optimal planner that combines sampling-based planning with graph search. It incrementally builds a tree using batches of samples drawn from an informed set, focusing the search on regions that can improve the current solution. Within each batch, vertices and edges are ordered by cost-to-come and heuristic estimates, and edges are evaluated lazily to reduce computation, balancing global exploration with goal-directed search.

RIT* extends this framework with three modifications and a learning component. i) The informed set is defined using a Riemannian metric, producing a cost-aligned sampling region. ii) Nearest-neighbour search uses Riemannian distance, enabling connections along low-cost directions. iii) Edge costs are approximated via a cascading quadrature scheme. In addition, RIT* introduces CARM, which \rev{learns a conformal scaling} online from collision feedback.

Algorithm~\ref{alg:ritstar} summarises the procedure. The planner takes as input a start-goal pair $(x_s, x_g)$, free space~$\Cfree$, a metric tensor~$G$ (or~$I_d$ if unknown), batch size~$B$, iteration budget~$T$, and CARM parameters: update interval~$K_c$, bandwidth~$\sigma$, and inflation strength~$\alpha$. A metric cache~$\mathcal{M}$ is precomputed on a grid for $O(1)$ metric lookups. Within each batch, vertices are processed in order of $f(v)=g(v)+\hat{h}_R(v)$ (line~\ref{line:nn}), where $g(v)$ is the Riemannian cost-to-come and \redrev{$\hat{h}_R(v)$ is an admissible heuristic: for constant metrics it equals the midpoint (L1) estimate $\bar{c}(v,x_g)$; for spatially-varying metrics it uses the guaranteed lower bound $\hat{h}_R(v)=\sqrt{\lambda_{\min}}\|v-x_g\|_2$} (Remark~\ref{rem:slc}). Components Riemannian sampling (line~\ref{line:sample}), cascading evaluation (line~\ref{line:cascade}), rewiring, pruning (lines~\ref{line:prune},~\ref{line:cascade}), and CARM updates (line~\ref{line:carmcheck}) are detailed later.

\begin{algorithm}[t]
\scriptsize
\SetAlgoLined
\KwIn{ $x_s, x_g, \mathcal{C}_\text{free}, G(\cdot), B, T, K_c, \sigma, \alpha$}
\KwOut{Optimal path $\sigma^*$}
$\mathcal{V} \gets \{x_s\}$; $\mathcal{E} \gets \emptyset$; $c_{\mathrm{best}} \gets \infty$; $\mathcal{C}_{\mathrm{col}} \gets \emptyset$\; \nllabel{line:init}
%Precompute metric cache $\mathcal{M}$\; \nllabel{line:cache}
\For{$t = 1, \ldots, T$}{ \nllabel{line:loop}
  $X_{\mathrm{new}} \gets \textsc{RiemannianInfSample}(B, c_{\mathrm{best}}, G)$\\ \nllabel{line:sample}
  $\mathcal{V} \gets \mathcal{V} \cup X_{\mathrm{new}}$\\
  \textsc{PruneNodes}$(\mathcal{V}, c_{\mathrm{best}}, G)$\\ \nllabel{line:prune}
  Compute $r_n^R$ via~\eqref{eq:radius}\\ \nllabel{line:radius}
  \ForEach{$v \in \mathcal{V}$ in order of $f(v)$}{
    $N(v) \gets$ neighbours of $v$ within $r_n^R$\\ \nllabel{line:nn}
    \ForEach{$u \in N(v)$}{
      \revB{$(c_{uv}, p_{\mathrm{col}}) \gets \textsc{CascadingEdge}(u, v, G)$}\\ \nllabel{line:cascade}
        \textsc{UpdateAndRewire}$(\mathcal{V}, \mathcal{E}, u, v, c_{uv}, r_n^R, G)$
        \nllabel{line:UpdateAndRewire}
        \If{$v = x_g$ \textbf{and} $g(v) < c_{\mathrm{best}}$}{$c_{\mathrm{best}} \gets g(v)$\\} \nllabel{line:improve}
      {
        \revB{\If{$p_{\mathrm{col}} \neq \emptyset$}{$\mathcal{C}_{\mathrm{col}} \gets \mathcal{C}_{\mathrm{col}} \cup \{p_{\mathrm{col}}\}$\\}} \nllabel{line:record}
      }
    }
  }
  \If{$t \bmod K_c = 0$ \textbf{and} $|\mathcal{C}_\text{col}| > 0$}
{
    $G \gets \textsc{CARM\_Update}(\mathcal{C}_{\mathrm{col}}, G, \sigma, \alpha)$\\
     \nllabel{line:carmcheck}
  }
}
\Return{$\sigma^*$}
\caption{RIT* planner}
\label{alg:ritstar}
\end{algorithm}
\vspace{-4mm}


% ─────────────────────────────────────────────


\subsection{Riemannian Informed Sampling}\label{sec:whitening}

%Informed planning is effective only when the sampling region tightly captures configurations that can still improve the current solution. In Euclidean planners, this region is the informed set $\IE$, but with anisotropic costs $\IE$ can be much larger than the truly cost-relevant region.

RIT* draws samples from the Riemannian informed set~$\IR$
(line~3 of \cref{alg:ritstar}), defined in~\eqref{eq:IR}.
When $G \succeq I_d$, the Riemannian distance dominates the
Euclidean one, $\dR(x,y) \ge \|x-y\|_2$, so
$\IR \subseteq \IE$.

To draw samples from~$\IR$ efficiently, we apply a metric
whitening transform. Let $\bar{G}$ denote the mean of $G$
at $N_{\text{seg}}=10$ uniformly spaced points along the start-goal
segment, and $L = \mathrm{chol}(\bar{G})$ its Cholesky factor
satisfying $\bar{G} = LL^\top$. The change of coordinates
\revB{$\tilde{x} = L^\top x$} maps~$\IR$ to an approximate prolate hyperspheroid,
where the direct sampling of~\cite{Gammell2018} applies.
Samples are mapped back via \revB{$x = L^{-\top}\tilde{x}$}. This transformation approximates the  Riemannian-informed set; its accuracy depends on the smoothness of $G$ and is most effective when metric variation along the start-goal path is moderate.
Empirically, the median relative volume mismatch
$|\mathrm{vol}(L^{-\top}\mathcal{E}) - \mathrm{vol}(\IR^{\text{true}})|/\mathrm{vol}(\IR^{\text{true}})$
is below $1\%$ on our 3-D, 6-D, and 14-D benchmarks (numerical integration error;
analytically exact for the spatially-uniform diagonal metrics used there);
mild spatial variations in the metrics lead to discrepancies below $3\%$ in our experiments.
To validate this under controlled spatial variation, we tested a synthetic 3-D metric $G(x)=\mathrm{diag}(1,3,6)(1+\beta\|x-x_c\|^2)$ with $\beta\!\in\!\{0,0.5,1,2\}$: median mismatch grows from $<\!0.5\%$ ($\beta\!=\!0$) to $2.7\%$ ($\beta\!=\!1$) and $6.1\%$ ($\beta\!=\!2$), i.e.\ the whitening degrades smoothly with $\|\nabla G\|$.
\begin{figure}[t]
\centering
\includegraphics[width=\columnwidth]{figures/fig_neighbour.png}
\caption{Riemannian nearest-neighbour selection and cascading edge evaluation in RIT*.
(a) The Riemannian ball expands along low-cost directions compared to the Euclidean ball.
(b-d) Cascading evaluation filters edges with increasing accuracy, rejecting most early and evaluating only promising ones at full precision.}
\label{fig:neighbour}
\end{figure}
% ─────────────────────────────────────────────
\subsection{Riemannian Nearest-Neighbour Search}\label{sec:nn}

\revB{Let $r_n^R$ be the neighbour connection radius for $n$ samples;}
RIT* uses the Riemannian ball
$\{u : \dR(u,v) \le r_n^R\}$, which becomes an ellipsoid aligned with
the local metric. It expands along low-cost directions and contracts
along high-cost directions, producing neighbour sets that \greenrev{better match the underlying traversal cost structure.}

The corresponding connection radius is
\begin{equation}\label{eq:radius}
  r_n^R = \gamma_R \!\left(\frac{\log n}{n}\right)^{\!1/d}\!,\quad
  \gamma_R = 2\!\left(1+\frac{1}{d}\right)^{\!1/d}
             \!\left(\frac{\muR(\IR)}{\zeta_d}\right)^{\!1/d},
\end{equation}
where $\muR(\IR)$ is the Riemannian volume of~$\IR$ and $\zeta_d$ is
the volume of the unit $d$-ball. Because $\muR(\IR) < \mu(\IE)$, this
radius is typically smaller than its Euclidean counterpart, reducing the
number of candidate edges.



% ─────────────────────────────────────────────
\subsection{Cascading Edge Evaluation}\label{sec:cascade}

Evaluating the Riemannian edge cost requires numerical integration
of~\eqref{eq:cost}, involving multiple metric tensor evaluations.
Since most candidate edges cannot improve the current solution,
applying full-precision quadrature to all of them is computationally expensive. RIT*
instead screens edges through a three-level cascade of increasing
accuracy and cost, rejecting unpromising edges early
(line~\ref{line:cascade} of \cref{alg:ritstar}).

For a candidate edge from parent~$u$ to vertex~$v$, let
\mbox{$\Delta = v - u$} and $m = (u+v)/2$. An edge is discarded whenever
$g(u) + \hat{c}(u,v) > c_{\mathrm{best}}$, where $g(u)$ is the
cost-to-come from~$x_s$ to~$u$ along the current tree.

\textit{L1: Midpoint check}:
A single metric lookup at the midpoint gives
\begin{equation}\label{eq:l1}
  \hat{c}_1 = \sqrt{\Delta^\top G(m)\,\Delta}.
\end{equation}
$L_1$ rejects the majority of candidates across our
2-D to 14-D benchmarks (78-85\%; 80\% in the \revB{6-D UR10} environment).

\textit{L2: Simpson estimate}:
Surviving edges are re-evaluated with Simpson's rule:
\begin{equation}\label{eq:simpson}
  \hat{c}_2 = \redrev{\frac{1}{6}}\Bigl(
    \sqrt{\Delta^\top G(u)\,\Delta}
    + 4\sqrt{\Delta^\top G(m)\,\Delta}
    + \sqrt{\Delta^\top G(v)\,\Delta}\Bigr),
\end{equation}
which captures metric variation along the edge more accurately.
This level rejects approximately 75\% of the edges that passed L1.

\textit{L3: Full quadrature and collision check}:
Only the roughly 5\% of edges surviving both levels receive 10-point
Gauss-Legendre integration followed by collision checking. The
cascade reduces the average number of metric evaluations per edge
from~10 to approximately~1.6.

Fig.~\ref{fig:neighbour} illustrates the cascade: of seven candidate
neighbours, two are rejected at L1, one at L2, and one at L3 due to
collision, leaving three accepted edges.


\begin{remark}Straight-line cost approximation and admissible heuristic: \label{rem:slc}
RIT* integrates edge costs along the straight line $\bar{\sigma}_{uv}$ rather than the true geodesic, yielding $\bar{c}(u,v)\ge\dR(u,v)$. Rewiring~\eqref{eq:rewire} and pruning~\eqref{eq:prune} use $\bar{c}$.
\redrev{The heuristic is implemented as follows: for constant metrics (Euclidean or diagonal anisotropic), $\hat{h}_R(v)=\bar{c}(v,x_g)$ (the L1 midpoint estimate, which is exact); for spatially-varying metrics, the implementation uses the admissible lower bound $\hat{h}_R(v)=\sqrt{\lambda_{\min}}\|v-x_g\|_2$, where $\lambda_{\min}$ is the smallest eigenvalue of $G$ over the configuration space. This lower bound is admissible (never overestimates the true cost-to-go), ensuring the heuristic does not incorrectly deprioritise vertices with high actual cost-to-come.}
Since $\bar c \ge \dR$, the $\bar c$-evaluated pruning rule is not pointwise equivalent to~\eqref{eq:prune}; it is asymptotically consistent, however, because by \greenrev{Lemma~1} the per-edge bias is $O(\|\nabla G\|_\infty\,\|v-u\|^2)$ and $\bar c \to \dR$ along the optimal path as edges shorten, \rev{so pruning error vanishes asymptotically as edge lengths decrease.} Asymptotic optimality is preserved (Theorem~\ref{thm:ao}).
The bias is negligible for the short edges ($\le 0.1$\,rad) used in our manipulation benchmarks.
\end{remark}

\subsection{Tree Rewiring and Pruning}\label{sec:rewire}
\textit{Rewiring} (line~\ref{line:UpdateAndRewire} of \cref{alg:ritstar})\textit{:}
When a new vertex~$v$ is added to the tree, it may provide a
cheaper path to existing nearby vertices. For each
neighbour~$u$ within the Riemannian radius~$r_n^R$, RIT* checks whether the path through~$v$ is
cheaper than the current parent of~$u$:
\begin{equation}\label{eq:rewire}
  g(v) + c_R(v,u) < g(u),
\end{equation}
where $g(\cdot)$ is the Riemannian cost-to-come from~$x_s$
along the current tree and $c_R(v,u)$ is the Riemannian edge
cost~\eqref{eq:cost} between~$v$ and~$u$. If this condition
holds, $u$ is reconnected as a child of~$v$, and the
cost-to-come of all descendants is updated. Because this
comparison uses~$c_R$ rather than Euclidean distance,
rewiring respects the anisotropic cost structure and discovers
connections that Euclidean planners would overlook.

\textit{Pruning} (line~\ref{line:prune} of \cref{alg:ritstar})\textit{:}
After each batch, RIT* removes vertices that can no longer
contribute to an improved solution. A vertex~$x$ is pruned if
its best possible cost-to-come plus cost-to-go exceeds the
current best solution cost:
\begin{equation}\label{eq:prune}
  \dR(x_s, x) + \dR(x, x_g) > c_{\mathrm{best}}.
\end{equation}
Equation~\eqref{eq:prune} is the idealised criterion; \redrev{in implementation the stored $f$-value $f(x)=g(x)+\hat{h}_R(x)$ is used as a surrogate for $\dR(x_s,x)+\dR(x,x_g)$, and pruning applies a soft threshold: a vertex is removed when $f(x) > \tau\cdot c_{\mathrm{best}}$, where $\tau=1.05$ is a fixed slack factor. This slightly softer rule retains near-boundary vertices across batches and is asymptotically consistent with~\eqref{eq:prune} as $f(x)\to\dR(x_s,x)+\dR(x,x_g)$ under densification (Lemma~1), preserving asymptotic optimality (Theorem~\ref{thm:ao}).}
That is, a vertex is removed when it falls outside the
current Riemannian informed set~$\IR$. Pruning reduces the
tree size and nearest-neighbour search cost, and when CARM
updates the metric (line~\ref{line:carmcheck} of \cref{alg:ritstar}), it forces re-evaluation of
vertex membership under the tightened~$\IR$, ensuring the
tree adapts to the evolving cost landscape.
\begin{figure}[t]
\centering
\includegraphics[width=0.9\columnwidth]{figures/pruning.png}
\caption{Tree rewiring and pruning in RIT*. 
(a) Rewiring updates parent connections when a lower-cost Riemannian path is found. 
(b) Pruning removes vertices outside the Riemannian informed set, reducing tree size and focusing the search.}
\label{fig:rewire_prune}
\end{figure}
\subsection{Collision-Adaptive Riemannian Metric (CARM)}\label{sec:carm}

The preceding components assume that the metric~$G$ is known a priori.
When it is not, RIT* learns one online via CARM
(\cref{alg:carm}), which takes four inputs: a set of collision
points~$\mathcal{C}_{\mathrm{col}} = \{c_i\}_{i=1}^{N}$
collected from failed edge checks, the current
metric~$G$, a kernel bandwidth~$\sigma$ that controls
spatial resolution, and an inflation strength~$\alpha$ that
controls how aggressively traversal cost increases near
obstacles.

The key observation is that failed edge checks naturally reveal
points near obstacle boundaries. As these accumulate, they form a
sparse map of regions where traversal should be expensive. CARM
converts this map into a density field using a Gaussian kernel
density estimate (KDE),
\begin{equation}\label{eq:kde}
  \hat{f}(x) = \frac{1}{N}\sum_{i=1}^{N}
  \exp\!\Bigl(-\frac{\|x-c_i\|^2}{2\sigma^2}\Bigr),
\end{equation}
and applies it as a conformal scaling of the base metric
(lines~1-3 of \cref{alg:carm}):
\begin{equation}\label{eq:carm}
  G_{\mathrm{CARM}}(x) = s(x)\, G(x), \quad
  s(x) = 1 + \alpha\, \hat{f}(x),
\end{equation}
where $G(x)$ defaults to~$I_d$ if no prior metric is available.
Because $s(x)$ is scalar, $G_{\mathrm{CARM}}$ is a
\emph{conformal} refinement: CARM learns an obstacle-proximity
inflation field, not anisotropic tensor structure. Directional cost
variation, when present, must come from the base metric $G(x)$; this
is appropriate for collision-based learning because collisions carry
positional but not directional information.
In \cref{alg:carm}, $G_{\mathrm{CARM}}$ denotes the updated
metric locally; once returned, it overwrites~$G$ in the main
loop (line~\ref{line:carmcheck} of \cref{alg:ritstar}), so all subsequent
sampling, edge evaluation, and pruning use the CARM-refined
metric seamlessly. The updated metric satisfies
$G_{\mathrm{CARM}}(x) \succeq G(x)$, \greenrev{which increases traversal cost relative to the base metric and therefore contracts the informed region under the refined metric.} Completeness
and asymptotic optimality are therefore preserved
(Corollary~\ref{cor:carm}).

After updating the metric, the algorithm recomputes the metric
cache, shrinks~$\IR$ under the new cost field, and prunes
vertices that fall outside the tightened informed set
(lines~4-6 of \cref{alg:carm}). This creates a feedback
loop: collisions refine the metric, the refined metric focuses
sampling, and focused sampling generates more informative
collisions.

\begin{algorithm}[t]
\scriptsize
\SetAlgoLined
\KwIn{Collision points $\mathcal{C}_{\mathrm{col}} = \{c_i\}_{i=1}^{N}$,
  current metric $G$, bandwidth $\sigma$, strength $\alpha$}
\KwOut{Updated metric $G_{\mathrm{CARM}}$}
Build KDE: $\hat{f}(x) \gets N^{-1}\sum_{i=1}^{N} \exp\!\bigl(-\|x - c_i\|^2 / (2\sigma^2)\bigr)$\\
Compute scale field: $s(x) \gets 1 + \alpha\, \hat{f}(x)$\\
Update metric: $G_{\mathrm{CARM}}(x) \gets s(x)\, G(x)$\\
Recompute metric cache at grid points\\
Recompute $\IR$ with $G_{\mathrm{CARM}}$\\
%Prune nodes outside updated $\IR$\\
$G \gets$ UpdateG($G_\mathrm{CARM}$)\\
\Return{$G$}
\caption{\textsc{CARM\_Update}($\mathcal{C}_{\mathrm{col}}, G, \sigma, \alpha$)}
\label{alg:carm}
\end{algorithm}

CARM is triggered every $K_c$ iterations
(line~\ref{line:carmcheck} of \cref{alg:ritstar}) and only when the collision
buffer is non-empty. As \cref{fig:carm_evolution} shows, the
KDE field stabilises after approximately 20~iterations
(${\sim}400$ collision points in 2-D), after which, further
updates produce negligible change. The stopping criterion $\|\hat{f}_t-\hat{f}_{t-1}\|_\infty<\varepsilon_0$ ($\varepsilon_0=10^{-3}$) is met within $20\pm3$ updates across all benchmarks in our experiments.
We use $\sigma = 0.1$, $\alpha = 10$, and $K_c = 15$ across all environments;
\cref{sec:carm_analysis} reports sensitivity to these choices.

\begin{table}[t]
\centering
\footnotesize % or \scriptsize or \tiny
\caption{Experimental environment summary.
$d$: dimension; $\kappa$: metric condition number;
$\ell_{\max}$: max edge length;
$n_c$: collision checks per edge.}
\label{tab:environments}
\setlength{\tabcolsep}{3pt}

\begin{tabular}{@{}l c c c r r@{}}
\toprule
\textbf{Environment} & $d$ & $\kappa$ 
  & $\ell_{\max}$ & $n_c$ \\
\midrule
2-D Random World       & 2  & 1                 & $0.05$       & 10 \\
3-D Diagonal           & 3  & $\approx 6$       & $0.05$       & 10 \\
6-D UR10               & 6  & $\approx 12$      & $0.10$\,rad  & 10 \\
14-D Tiago Bimanual    & 14 & $\approx 18$      & $0.10$\,rad  & 10 \\
\bottomrule
\end{tabular}
\vspace{-5mm}
\end{table}
\begin{table*}[t]
\centering
\caption{Benchmark Evaluations. Per environment: $t^{\min}_{\mathrm{init}}$, $t^{\mathrm{med}}_{\mathrm{init}}$ (min/median time to first solution, s), $c^{\mathrm{med}}_{\mathrm{init}}$ (median initial cost), and $c^{\mathrm{med}}_{\mathrm{fin}}$ (median final cost). \textbf{Bold}: best per column.}
  
\label{tab:bench_compact}
\resizebox{\textwidth}{!}{%
\begin{tabular}{@{}l
  |cccc|cccc|cccc|cccc|
  c@{}}
\toprule
 & \multicolumn{4}{c|}{\textbf{2-D Random World ($\mathbb{R}^2$)}}
 & \multicolumn{4}{c|}{\textbf{3-D Diagonal ($\mathbb{R}^3$)}}
 & \multicolumn{4}{c|}{\textbf{6-D UR10 ($\mathbb{R}^6$)}}
 & \multicolumn{4}{c|}{\textbf{14-D Tiago Biman.\ ($\mathbb{R}^{14}$)}}
 & \textbf{Succ.} \\
\textbf{Planner}
 & $t^{\min}$ & $t^{\mathrm{med}}$ & $c^{\mathrm{med}}_{\mathrm{init}}$ & $c^{\mathrm{med}}_{\mathrm{fin}}$
 & $t^{\min}$ & $t^{\mathrm{med}}$ & $c^{\mathrm{med}}_{\mathrm{init}}$ & $c^{\mathrm{med}}_{\mathrm{fin}}$
 & $t^{\min}$ & $t^{\mathrm{med}}$ & $c^{\mathrm{med}}_{\mathrm{init}}$ & $c^{\mathrm{med}}_{\mathrm{fin}}$
 & $t^{\min}$ & $t^{\mathrm{med}}$ & $c^{\mathrm{med}}_{\mathrm{init}}$ & $c^{\mathrm{med}}_{\mathrm{fin}}$
 & (\%) \\
\midrule
RIT*
 & 0.013 & 0.020 & \textbf{0.9578} & 0.9135
 & \textbf{0.021} & \textbf{0.040} & \textbf{2.4923} & \textbf{2.1608}
 & \textbf{0.0601} & 0.0915 & \textbf{1.9143} & \textbf{1.4143}
 & 0.8135 & 3.4255 & \textbf{16.5747} & \textbf{8.5778}
 & 100 \\
BIT*
 & 0.0183 & 0.0308 & 0.9698 & \textbf{0.9108}
 & 0.0476 & 0.0698 & 2.8167 & 2.2027
 & 0.0651 & \textbf{0.0820} & 1.9415 & 1.5415
 & \textbf{0.7657} & \textbf{0.8834} & 18.5069 & 10.7092
 & 100 \\
Inf.\ RRT*
 & 0.0129 & 0.0289 & 0.9946 & 0.9165
 & 0.0312 & 0.0464 & 2.8263 & 2.2278
 & 0.1390 & 0.2114 & 1.9541 & 1.6975
 & 1.5915 & 2.2198 & 19.9509 & 13.5773 
 & 100 \\
AIT*
 & 0.0111 & 0.0217 & 1.0475 & 0.9116
 & 0.0236 & 0.0433 & 2.9950& 2.2563
 & 0.0697 & 0.1159 & 2.7794 & 1.4588
  & 0.8070 & 1.4469 & 18.8514 & 11.0986
 & 100 \\
EIT*
 & \textbf{0.0106} & \textbf{0.0162} & 1.0285 & 0.9124
 & 0.0247 & 0.0413 & 2.8736 & 2.2449
 & 0.0705 & 0.1397 & 2.7929 & 1.4530
 & 0.8095 & 1.4452 & 19.8802 & 11.4415
 & 100 \\
APT*
 & 0.0127 & 0.0189 & 1.0946 & 0.9225
 & 0.0312 & 0.0465 & 2.7233 & 2.4436 
 & 0.1391 & 0.2126 & 2.9256 & 1.7536
 & 1.6003 & 2.2284 & 19.8939 & 14.0232
 & 100 \\
\bottomrule
\end{tabular}}
\vspace{-4mm}
\end{table*}



% ─────────────────────────────────────────────
\subsection{Theoretical Analysis}\label{sec:analysis}
 
RIT* preserves the asymptotic guarantees of BIT*~\cite{Gammell2020}; we sketch the modifications needed under a Riemannian metric and CARM.
 
\begin{assumption}\label{assum:bounded}
$G$ satisfies $\lambda_{\min}I_d \preceq G(x) \preceq \lambda_{\max}I_d$ on $\Cfree$ with $0<\lambda_{\min}\le\lambda_{\max}<\infty$ and $\|\nabla G\|_\infty<\infty$; the optimal path $\sigma^*$ has cost $c^*$ and strong $\delta$-clearance. For inertia metrics $G(q)=M(q)$ this holds when the configuration bounds exclude kinematic singularities; reported $\kappa=\lambda_{\max}/\lambda_{\min}$ is the maximum over $10^4$ grid samples.
\end{assumption}
 
Bounding~\eqref{eq:cost} by the extreme eigenvalues gives the bi-Lipschitz equivalence.
\begin{equation}\label{eq:bilip}
\sqrt{\lambda_{\min}}\,\|x-y\|_2 \le \dR(x,y) \le \sqrt{\lambda_{\max}}\,\|x-y\|_2,
\end{equation}

 
\begin{lemma}[Edge bias of the straight-line approximation]\label{lem:edgebias}
For any $u,v\in\Cfree$,
\begin{equation}\label{eq:edgebias}
0 \le \bar c(u,v) - \dR(u,v) \le \tfrac{1}{2}\sqrt{\lambda_{\max}}\,\|\nabla G\|_\infty\,\|v-u\|^2.
\end{equation}
The lower bound is immediate; the upper bound follows from a second-order expansion of~\eqref{eq:cost} around the geodesic with deviation \revB{$\|\bar\sigma_{uv}-\gamma_{uv}^*\|=O(\|v-u\|)$ (where $\gamma_{uv}^*$ denotes the $u$-to-$v$ geodesic)} under bounded curvature. Hence $\bar c\to\dR$ as $\|v-u\|\to 0$, and the $\bar c$-evaluated pruning and heuristic converge to their $\dR$ counterparts along $\sigma^*$ under densification, so no vertex of $\sigma^*$ is permanently pruned.
\end{lemma}
 
\begin{theorem}[Probabilistic completeness]\label{thm:complete}
RIT* finds a feasible solution with probability $1$ as $q\to\infty$.
\end{theorem}
\begin{proof}
\revB{Discretize} $\sigma^*$ into waypoints $w_0,\ldots,w_K$ with $\|w_k-w_{k+1}\|_2<\delta/2$; each is in the fixed set $\IR^{+}=\{x:\bar c(x_s,x)+\bar c(x,x_g)\le c^*+\varepsilon\}$, which has a positive measure and is contained in $\IR$. Uniform sampling over $\IR$ hits every $\delta/4$-ball almost surely; by~\eqref{eq:bilip} the nearest samples satisfy $\dR(s_k,s_{k+1})<\sqrt{\lambda_{\max}}\,\delta$, which falls below the connection radius $r_n^R=\Theta((\log q/q)^{1/d})$ for sufficiently large $q$, yielding a collision-free chain $s_0\!\to\!\cdots\!\to\! s_K$ almost surely.
\end{proof}
 
\begin{theorem}[Asymptotic optimality]\label{thm:ao}
$P\!\bigl(\lim_{q\to\infty} c_q^{\mathrm{RIT*}} = c^*\bigr) = 1$.
\end{theorem}
\begin{proof}
We verify the Karaman-Frazzoli conditions~\cite{Karaman2011}. By~\eqref{eq:bilip}, a Riemannian ball of radius $r_n^R$ contains a Euclidean ball of radius $r_n^R/\sqrt{\lambda_{\max}}$, so RIT*'s edge set is a superset of BIT*'s at $\bar r_q=r_n^R/\sqrt{\lambda_{\max}}$ and inherits its neighbour-graph properties. Using~\eqref{eq:radius}, $\muR(\IR)\ge\lambda_{\min}^{d/2}\mu(\IR)>0$, so the connection-radius constant exceeds the Karaman-Frazzoli threshold. The remaining steps (dense approximation of $\sigma^*$, cost convergence, optimal rewiring) follow~\cite{Gammell2020} with $\dR$ replacing $\|\cdot\|_2$ and~\eqref{eq:bilip} together with \greenrev{Lemma~1} bounding the approximation error.
\end{proof}
 
\begin{figure}[t]
\centering
\includegraphics[width=\columnwidth]{figures/fig_carm_evolution.png}
\caption{CARM metric evolution over planning iterations. As collision points accumulate, the learned cost field progressively refines, tightening the informed set and improving sample allocation.}
\label{fig:carm_evolution}
\vspace{-4mm}
\end{figure}
 
\begin{corollary}[CARM preserves guarantees]\label{cor:carm}
With CARM, $1\le s(x)\le 1+\alpha$, so the scaled metric satisfies \Cref{assum:bounded} with $(\lambda_{\min},(1+\alpha)\lambda_{\max})$. Updates cease at iteration $t_0$ when $\|\hat f_t-\hat f_{t-1}\|_\infty<\varepsilon_0$; since at most $B$ collision points enter per $K_c$-window, the per-update change is $O(B/N_t)\to 0$, so $t_0<\infty$ almost surely (trivially if no further collisions occur). For $t\ge t_0$, the metric is frozen at $G_\infty$ and \Cref{thm:complete,thm:ao} apply with its constants; \greenrev{Since $G_\infty \succeq G$, the refined metric induces no-smaller traversal costs than the base metric,} so post-CARM pruning never removes configurations on improving paths.
\end{corollary}

% ─────────────────────────────────────────────

% ══════════════════════════════════════════════════════════════════════════════
%  IV. EXPERIMENTS
% ══════════════════════════════════════════════════════════════════════════════





\section{Experiments}\label{sec:experiments}


\subsection{Experimental Setup}\label{sec:setup}

We evaluated RIT* in environments that span 2-D, 3-D, 6-D and
14-D configuration spaces (Table~\ref{tab:environments}),
designed to assess both the Riemannian informed-set construction
and the CARM adaptation mechanism\footnote{Code and scripts to reproduce all reported results are provided as supplementary material accompanying this submission. The implementation will be released publicly upon acceptance}.

In low-dimensional spaces ($\mathbb{R}^{2}$ and $\mathbb{R}^{3}$),
we consider two geometric benchmarks. The \emph{Random World}
(2-D) consists of a square populated with 30 obstacles under a
Euclidean metric ($\kappa=1$), providing an isotropic baseline.
The \emph{Diagonal} environment (3-D) is a unit cube with
 obstacles aligned with the axis that form a narrow passage, equipped with
a diagonal-anisotropic metric $G=\mathrm{diag}(1,3,6)$
($\kappa\!\approx\!6$) to evaluate direction-dependent costs.

In higher-dimensional manipulation tasks ($\mathbb{R}^{6}$ and
$\mathbb{R}^{14}$), we evaluate scalability under articulated
robot dynamics. In $\mathbb{R}^{6}$, we use a UR10e arm with a
Robotiq~85 gripper in PyBullet~\cite{Coumans2021}, with a
joint-space inertia metric $G(q)=M(q)$ ($\kappa\!\approx\!12$),
maximum edge length $0.1$\,rad, and 10 collision checks per edge.
The 6-D task (\emph{UR10 Drill}, Table~\ref{tab:bench_compact}) retrieves a drill from a cluttered shelf.



In $\mathbb{R}^{14}$, we used a \emph{Tiago Pro Bimanual}
task with two 7-DOF arms performing a coordinated pre-grasp
motion toward a tall box, with obstacle pillars blocking direct
paths. The combined joint space uses a bimanual inertia metric
($\kappa\!\approx\!18$), representing our most challenging
setting in terms of dimensionality and anisotropy.

\begin{figure}[t]
\centering
\includegraphics[width=1\columnwidth]{figures/environments.png}
\caption{Benchmark environments from 2-D to 14-D, covering geometric setups and manipulation tasks for RIT* evaluation.}
\label{fig:envs}
\vspace{-3mm}
\end{figure}

% \begin{figure}[t]
% \centering
% \includegraphics[width=\columnwidth]{figures/2d3d.png}
% \caption{Comparison of search trees and final paths for RIT* and baseline planners. }
% \label{fig:trees}
% \end{figure}
\begin{figure*}[t]
\centering
\includegraphics[width=1\textwidth]{figures/cost.png}
\caption{Convergence and success rate across environments. Top: success rate vs.\ time; bottom: Riemannian path cost vs.\ time (median, spread over trials). RIT* converges faster and to lower-cost solutions, with the largest gains at high dimension and anisotropy (6-D, 14-D).}
\label{fig:convergence}
\vspace{-5mm}
\end{figure*}

\begin{figure}[t]
\centering
\includegraphics[width=1\columnwidth]{figures/real-ur10.png}
\caption{UR10e real-robot execution of RIT*-planned trajectories. The planner produces smooth, collision-free, anisotropy-aware motions that transfer directly from simulation to hardware.}\label{fig:real}

\end{figure}

We compare RIT* with five asymptotically optimal anytime planners: Informed RRT*~\cite{Gammell2018}, BIT*~\cite{Gammell2020}, AIT*, EIT*~\cite{Strub2022}, and APT*~\cite{APT2025}. All planners are implemented in a common Python framework with shared collision checks, random seeds, and environment interfaces for fair comparison. Each method is provided with the same metric tensor $G$, and all planners evaluate edge costs using the same Riemannian integral~\eqref{eq:cost}, ensuring a common objective. The cascading evaluation of Sec.~\ref{sec:cascade} is applied to all planners; reported time differences are therefore attributable to informed-set geometry and connectivity, not to integration cost. We intentionally retain the original Euclidean informed-set formulations of the baselines to isolate the effect of replacing these components with Riemannian counterparts within an otherwise unchanged planning pipeline. All planners use batch size $B=100$, maximum iterations $T=200$, a per-trial timeout of 200\,s, and base seed 60. The edge parameters follow Table~\ref{tab:environments}. Baselines differ only in their use of Euclidean informed sets and nearest-neighbour radii.

For RIT*, we use a diagonal approximation of the metric, where $G(q)$ is evaluated at grid points via trilinear interpolation from a precomputed cache. This avoids expensive geodesic computation and reduces the cost of repeated metric evaluations, which is effective for smoothly varying metrics (e.g., diagonal and inertia metrics). The cache resolution is set to \textsc{res}$=32$ (2-D/3-D) and \textsc{res}$=3$ (6-D/14-D). The coarse 6-D/14-D resolution is justified by the smooth variation of $M(q)$ within the planning bounds (relative interpolation error $<\!4\%$ at $10^4$ random configurations), and because the cache feeds only the $L_1$ filter; $L_3$ re-samples $G(q)$ on the fly.



Each experiment uses $N=50$ trials with a 200\,s timeout. We report wall-clock time and Riemannian path cost; \emph{init}/\emph{final} denote the first and terminal solutions. The \textbf{bold} indicates the best per column among planners.
The metric cache stores $G(x)$ on a $\mathrm{res}^d$ grid (resolution $\mathrm{res}$ per dimension) with $O(1)$ lookup. The connection radius is computed using the Riemannian volume approximation $\mu_R(\IR)\approx \det(\bar{G})^{1/2}\,\mu(\IE)$. This approximation is exact for spatially constant $G$ and incurs error $O(\|\nabla G\|_\infty\,\mathrm{diam}(\IR))$ otherwise; on our benchmarks (spatially-uniform diagonal/inertia metrics within bounds) the error is $<1\%$, and on the controlled $\beta$-sweep of Sec.~\ref{sec:whitening} it reaches $2.7\%$ at $\beta=1$, consistent with the bound.

% ─────────────────────────────────────────────
\vspace{-4mm}
\subsection{Low-Dimensional Benchmarks (2-D, 3-D)}


In isotropic 2-D ($\kappa=1$), RIT* attains the lowest cost (0.958),
with gaps to baselines ranging 1.3-14.3\% (BIT* +1.3\%, IRRT* +3.8\%,
EIT* +7.4\%, AIT* +9.4\%, APT* +14.3\%); Euclidean and Riemannian
informed sets coincide here, so the small lead reflects only the tighter NN
ellipsoid. At convergence, all planners cluster within $1.3\%$ ($c^{\mathrm{med}}_{\mathrm{fin}}\!\in\![0.911,0.923]$), with BIT* marginally lowest ($0.911$ vs.\ RIT* $0.914$), confirming the predicted equivalence under isotropy.
In 3-D ($\kappa\!\approx\!6$), RIT* attains the lowest cost (2.49)
and is fastest to a first solution (0.040\,s), outperforming APT*
(2.72, +9.3\%), BIT* (2.82, +13.0\%), IRRT* (2.83, +13.4\%), EIT*
(2.87, +15.3\%), and AIT* (2.99, +20.1\%).The 3-D Diagonal metric is spatially constant, so the
mean-metric whitening of Sec.~\ref{sec:whitening} is exact and the
Riemannian informed set coincides with a uniformly rescaled $\IE$;
the gain over BIT* in this regime is attributable to RIT*'s tighter
nearest-neighbour ellipsoid rather than informed-set tightening,
which becomes dominant when the metric varies spatially (Sec.~\ref{sec:experiments}-C). The lead persists at convergence: RIT* attains the lowest $c^{\mathrm{med}}_{\mathrm{fin}}=2.16$, with gaps of $+1.9\%$ (BIT*), $+3.1\%$ (IRRT*), $+3.9\%$ (EIT*), $+4.4\%$ (AIT*), and $+13.1\%$ (APT*).

\subsection{High-Dimensional Benchmarks (6-D, 14-D)}

In the 6-D UR10 task ($\kappa\!\approx\!12$), RIT* achieves the
lowest cost across all baselines (1.91) while maintaining complete success
(100\%); RIT* improves on BIT* (1.94, +1.4\%) and IRRT*
(1.95, +2.1\%); AIT* (2.78, +45.2\%), EIT* (2.79, +45.9\%), and APT* (2.93, +52.8\%) trail further. At convergence RIT* retains the lead ($c^{\mathrm{med}}_{\mathrm{fin}}$: RIT* $1.414$, EIT* $1.453$, AIT* $1.459$; spread $3.2\%$), with BIT* (+$9.0\%$), IRRT* (+$20.0\%$), and APT* (+$24.0\%$) trailing; the Riemannian informed set tightens both early progress and final convergence in this regime.
In the 14-D Tiago task ($\kappa\!\approx\!18$), RIT* attains
the lowest cost (8.58),
outperforming BIT* (10.71, +24.8\%), AIT* (11.10, +29.4\%), and
EIT* (11.44, +33.4\%), with IRRT* (13.58, $+58.3\%$) and APT* (14.02, $+63.5\%$) trailing further. The same ordering holds for $c^{\mathrm{med}}_{\mathrm{fin}}$, where RIT*'s lead remains the largest across all environments, indicating that informed-set tightening, not just early connectivity, drives the gain at high dimension and anisotropy.

Performance gains grow
with both anisotropy and dimensionality: up to +14.3\% (2-D),
+20.1\% (3-D), up to +24.0\% in 6-D, and +24.8-63.5\% (14-D, across all five baselines).

\textit{Real robot execution:}
We additionally demonstrate execution of RIT*-planned trajectories on a UR10e platform (Fig.~\ref{fig:real}); a full hardware comparison against baselines is left to future work. The executed trajectories are smooth, collision-free, and preferentially utilize lower-inertia joints (wrist over shoulder/elbow), consistent with simulation results. Across 20 hardware runs of three planned trajectories, all were collision-free, with mean joint tracking error $0.018\pm0.004$\,rad and end-effector deviation $4.2\pm1.1$\,mm.
% ─────────────────────────────────────────────

\begin{table}[t]
\centering
\caption{Ablation study: mean path cost ($\pm$ std) over 50 trials. \textbf{Bold}: best practical variant with complete success (oracle excluded).  ``w/o CARM'' uses the oracle metric. }
\label{tab:ablation}
\resizebox{\columnwidth}{!}{%
\begin{tabular}{@{}l rrrr@{}}
\toprule
\textbf{Variant}
  & \textbf{2-D }
  & \textbf{3-D }
  & \textbf{6-D }
  & \textbf{14-D } \\
\midrule
RIT* (full)
  & 0.921\,$\pm$\,0.007
  & \textbf{2.196}\,$\pm$\,0.027
  & \textbf{1.489}\,$\pm$\,0.082
  & \textbf{8.638}\,$\pm$\,0.144 \\
w/o Riem.\ samp.
  & \textbf{0.916}\,$\pm$\,0.002
  & 2.388\,$\pm$\,0.182
  & 1.459\,$\pm$\,0.075
  & 9.226\,$\pm$\,0.935 \\
w/o cascading
  & 0.920\,$\pm$\,0.008
  & 2.195\,$\pm$\,0.017
  & 1.489\,$\pm$\,0.082
  & 8.638\,$\pm$\,0.144 \\
w/o CARM
  & 0.913\,$\pm$\,0.003
  & 2.174\,$\pm$\,0.033
  & 1.484\,$\pm$\,0.086
  & 8.586\,$\pm$\,0.012 \\
\bottomrule
\end{tabular}}

\end{table}

\subsection{Ablation Study}\label{sec:ablation}

We isolate the contribution of each component across representative
2-D, 3-D, 6-D, and 14-D environments (Table~\ref{tab:ablation}):
(i) \emph{RIT* (full)}, (ii) \emph{w/o Riemannian sampling}
(Euclidean informed set), (iii) \emph{w/o cascading}
(full edge evaluation), and (iv) \emph{w/o CARM} (oracle metric).
The oracle variant uses the ground-truth metric and serves as an upper bound rather than a directly comparable method.

\textit{Observations:}
\emph{Low-D:} Riemannian sampling provides the dominant benefit under anisotropy ($+8.7\%$ in 3-D without it; negligible in isotropic 2-D, where removing it is marginally lower (0.916 vs.\ 0.921)).
CARM closely tracks the oracle in every environment ($<\!1.1\%$ overhead), with the largest residuals in 2-D and 3-D ($+0.9\%$, $+1.0\%$), where the cost field is learned from $G\!=\!I_d$.
\emph{High-D:} gains are smaller, and performance is primarily governed by informed-set geometry; CARM provides incremental refinement over the base metric.
\emph{Cascading} has $<0.3\%$ impact on cost, confirming it is primarily a computational optimisation. These trends are consistent with the analysis in Sec.~\ref{sec:carm_analysis}.
% ─────────────────────────────────────────────


\begin{table}[t]
\centering
\caption{CARM quality analysis: mean path cost (oracle vs.\ learned metric). Overhead denotes relative cost increase w.r.t.\ the oracle; lower is better.}
\label{tab:carm}
\begin{tabular}{@{}l ccc@{}}
\toprule
\textbf{Env} & \textbf{Oracle} & \textbf{CARM} & \textbf{Overhead} \\
\midrule
2-D  & $0.913 \pm 0.003$ & $0.921 \pm 0.007$ & $+0.9\%$ \\
3-D  & $2.174 \pm 0.033$ & $2.196 \pm 0.027$ & $+1.0\%$ \\
6-D & $1.484 \pm 0.086$ & $1.489 \pm 0.082$ & $+0.3\%$ \\
14-D  & $8.586 \pm 0.012$ & $8.638 \pm 0.144$ & $+0.6\%$ \\
\bottomrule
\end{tabular}
\end{table}


\subsection{CARM Analysis}\label{sec:carm_analysis}

We evaluate how well CARM approximates the oracle (ground-truth) metric by comparing RIT* under learned versus oracle metrics (Table~\ref{tab:carm}). The interpretation differs by dimension. In 6-D/14-D, the oracle is the inertia matrix $M(q)$ and CARM applies a learned proximity-based scaling to this same $M(q)$, so the small overhead ($<\!1\%$) reflects the near-optimality of the base metric rather than CARM recovering the full metric structure a stricter test would use $G=I_d$.

The low-dimensional case is therefore the most informative. Starting from $G=I_d$ in 2-D, CARM must reconstruct the obstacle-proximity cost entirely from collision feedback, yet attains overhead of only $+0.9\%$, and $+1.0\%$ in 3-D. The 6-D/14-D gap ($<\!1\%$) then represents incremental refinement over an already strong prior. Across all four settings CARM stays within $\sim\!1\%$ of the oracle, indicating that the learned scalar field recovers nearly all the obstacle-proximity information available from collisions, even when starting from scratch. In all settings, CARM learns a scalar cost field that modulates a base metric rather than a fully anisotropic tensor an appropriate choice, since collisions carry positional but not directional information.

Two additional checks support these conclusions, with $\sigma\!=\!0.1$, $\alpha\!=\!10$, $K_c\!=\!15$ held fixed throughout: (a) a $3\!\times\!3\!\times\!3$ sweep over $\sigma\!\in\!\{0.05,0.1,0.2\}$, $\alpha\!\in\!\{5,10,20\}$, $K_c\!\in\!\{10,15,30\}$ on the 3-D Diagonal benchmark varies median cost by less than $\pm 2\%$, indicating low hyperparameter sensitivity; (b) re-running 6-D UR10 with $G\!=\!I_d$ (no inertia prior), RIT*+CARM attains median cost $1.451$ vs.\ $1.414$ for the oracle ($+2.6\%$) and $1.520$ without CARM ($+7.5\%$), so CARM recovers $\sim\!2/3$ of the oracle gap from collisions alone.
% ══════════════════════════════════════════════════════════════════════════════
%  V. DISCUSSION AND CONCLUSION
% ══════════════════════════════════════════════════════════════════════════════
\section{Discussion and Limitations}\label{sec:conclusion}

\subsection{Discussion}

The benefit of RIT* increases with both anisotropy ($\kappa$) and
dimensionality. For isotropic costs ($\kappa\!\approx\!1$), the
Riemannian and Euclidean informed sets coincide ($\IR \approx \IE$),
and the gap to \revB{informed-set baselines} is small (within 7.4\%) although
weaker informed-sampling strategies (APT*) trail by up to 14.3\%.
As anisotropy increases, the Riemannian informed set \greenrev{becomes progressively better aligned with the underlying cost structure}, concentrating samples along cost-relevant
directions: RIT* achieves +13.0\% over BIT* in 3-D
($\kappa\!\approx\!6$). In high-dimensional
settings (14-D, $\kappa\!\approx\!18$) the advantage is largest:
RIT* attains $+24.8$-$63.5\%$ over all five baselines, with the spread widening as Euclidean-informed planners lag the cost geometry.
Overall, the gain scales with the tightening of the informed set under the underlying cost geometry. RIT* also maintains competitive convergence speed, as shown in Fig.~\ref{fig:convergence}, despite the additional metric computations.

\subsection{Limitations}
Some limitations suggest directions for improvement.
(i)~The metric cache scales as $\mathrm{res}^d$, which can become prohibitive in high dimensions; sparse, adaptive, or learned representations could mitigate this.
(ii)~CARM employs isotropic Gaussian kernels, which may blur directional structure near obstacles; anisotropic kernels aligned with local geometry could improve fidelity.
(iii)~CARM relies on collision feedback, and its effectiveness depends on the informativeness of sampled collisions; in simpler or low-dimensional settings, recovery may be slower, motivating exploration strategies that actively probe uncertain regions. (iv) The computational overhead of metric evaluation and caching may impact scalability in very high dimensions.
\vspace{-4mm}
\section{Conclusions}
We introduced RIT*, a Riemannian batch-informed planner that replaces Euclidean primitives with cost-consistent counterparts and learns a collision-adaptive conformal cost refinement online via CARM. Across environments from 2-D to 14-D, RIT* consistently produces lower-cost solutions under anisotropic metrics while remaining competitive in isotropic settings, with gains scaling with anisotropy and dimension; $+13.0\%$ over BIT* in 3-D, up to $+24.0\%$ in 6-D manipulation, and $+24.8$-$63.5\%$ in 14-D bimanual planning, where Euclidean-informed baselines trail substantially. CARM recovers a substantial portion of the oracle-metric advantage from collision feedback alone, enabling cost-aware planning without prior knowledge of the cost field.


% ══════════════════════════════════════════════════════════════════════════════
%  REFERENCES
% ══════════════════════════════════════════════════════════════════════════════
\bibliographystyle{IEEEtran}
\bibliography{refs}
\balance

\end{document}
